\documentclass[Riemann_Hypothesis.tex]{subfiles}
\begin{document}


\pagebreak
\color{red}
\part{Introduction}
\color{black}

L'hypothèse de Riemann est l'un des problèmes ouverts les plus célèbres 
et les plus anciens des mathématiques. Énoncée il y a 167 ans, en 1859, par 
le mathématicien allemand Bernhard Riemann, elle demeure encore aujourd'hui 
au cœur des recherches en théorie analytique des nombres. 

Ce problème a été introduit dans l'unique article que Riemann a publié sur la théorie 
des nombres, intitulé \textit{Ueber die Anzahl der Primzahlen unter einer gegebenen 
Grösse} (Sur le nombre de nombres premiers inférieurs à une taille donnée). 
Dans ce mémoire, Riemann cherchait à comprendre la distribution des 
nombres premiers et à donner une expression exacte de la fonction de comptage des 
nombres premiers, souvent notée $\pi(x)$.

Pour ce faire, Riemann a étudié la fonction zêta, initialement introduite 
par Leonhard Euler pour les variables réelles.

\begin{DefNef}[label=df:FonctionZetaEuler]{Fonction zêta d'Euler-Riemann}
    Pour tout nombre complexe $s \in \mathbb{C}$ tel que $\Re(s) > 1$, la fonction zêta est définie par la série :
    \begin{equation*}
        \zeta(s) \;=\; \sum_{n=1}^{+\infty} \frac{1}{n^s}
    \end{equation*}
    Euler avait déjà démontré que cette série pouvait 
    s'exprimer sous la forme d'un produit eulérien portant sur l'ensemble des 
    nombres premiers $\mathbb{P}$ :
    \begin{equation*}
        \zeta(s) \;=\; \prod_{p \in \mathbb{P}} \left(1 - p^{-s}\right)^{-1}
    \end{equation*}
\end{DefNef}

La contribution majeure de Riemann a été d'étendre le domaine de définition de 
cette fonction à l'ensemble du plan complexe. Par le procédé de prolongement 
analytique, il a montré que la fonction $\zeta$ pouvait être définie pour tout 
$s \in \mathbb{C} \setminus \{1\}$, où elle admet un pôle simple. 

De plus, Riemann a mis en évidence une équation fonctionnelle reliant $\zeta(s)$ 
à $\zeta(1-s)$, révélant ainsi une symétrie fondamentale de la fonction. 
Il a identifié deux types de zéros pour cette fonction:
\begin{itemize}
    \item Les \textbf{zéros triviaux}, situés aux entiers pairs strictement négatifs ($-2, -4, -6, \dots$).
    \item Les \textbf{zéros non triviaux}, qui appartiennent tous à la bande critique définie par $0 < \Re(s) < 1$.
\end{itemize}

L'hypothèse de Riemann stipule que tous ces zéros non triviaux 
sont de partie réel égale à $\frac{1}{2}$.

\begin{theorem}[label=th:HypotheseRiemann]{Hypothèse de Riemann}
    Tous les zéros non triviaux de la fonction $\zeta$ ont pour partie réelle $\frac{1}{2}$. Autrement dit, pour tout $s \in \mathbb{C}$ tel que $0 < \Re(s) < 1$ et $\zeta(s) = 0$, on a :
    \begin{equation*}
        \Re(s) \;=\; \frac{1}{2}
    \end{equation*}
\end{theorem}

L'importance de cette conjecture réside dans son lien intime avec les nombres premiers. 
Riemann a en effet prouvé que la répartition des zéros non triviaux de la fonction zêta 
dicte les fluctuations de la distribution des nombres premiers. 

Dans le prolongement des travaux fondateurs de Riemann, le vingtième siècle a vu 
l'émergence de nouvelles approches pour aborder ce problème, notamment grâce aux 
contributions majeures du mathématicien français André Weil. 

Dans les années 1940, Weil s'est intéressé à un analogue de l'hypothèse de Riemann dans 
le cadre de la géométrie algébrique sur les corps finis. Il est parvenu à démontrer 
l'hypothèse de Riemann pour les courbes algébriques sur les corps finis. Ces résultats 
ont posé les jalons des célèbres conjectures de Weil, qui seront plus tard entièrement 
résolues par Pierre Deligne.

Outre cette avancée en géométrie algébrique, l'apport direct d'André Weil 
à la théorie analytique des nombres classique réside dans sa généralisation des formules 
explicites. Riemann avait initialement établi une formule reliant les nombres premiers 
aux zéros de la fonction zêta. En 1952, d'après ses 
travaux (voir \cite{weil1952_formules}), Weil a étendu 
cette relation à une vaste classe de fonctions tests régulières, 
révélant une dualité entre les puissances de nombres premiers 
et le spectre des zéros non triviaux.

Cette démarche l'a conduit à formuler un critère analytique équivalent 
à l'hypothèse de Riemann, connu aujourd'hui sous le nom de critère de 
positivité de Weil. Pour l'énoncer, on introduit une distribution spécifique, 
souvent appelée forme hermitienne de Weil, agissant sur un espace de fonctions 
tests approprié.

\begin{DefNef}[label=df:FormeHermitienneWeil_fr_intro]{Forme hermitienne de Weil}
    Soit $\mathcal{W}$ l'espace des fonctions tests de Weil sur le groupe 
    multiplicatif $\left(\mathbb{R}_+^*, \times\right)$. \par
    Soient $f$ et $g$ deux fonctions de $\mathcal{W}$. \par
    On pose $g^*$ la fonction involutive définie pour tout $x \in \mathbb{R}_+^*$ par :
    \begin{equation*}
        g^*(x) \;=\; \frac{1}{x} \, \overline{g\left(\frac{1}{x}\right)}
    \end{equation*}
    On note $\Phi = f * g^*$ le produit de convolution multiplicatif de $f$ et $g^*$. \par
    D'après les formules explicites d'André Weil, on appelle 
    \textbf{forme hermitienne de Weil}, que l'on notera $\langle f, g \rangle_{\mathcal{HP}}$, la forme sesquilinéaire définie par l'action de la distribution géométrique sur la fonction test $\Phi$ :
    \begin{equation*}
        \langle f, g \rangle_{\mathcal{HP}} \;:=\; W_{\text{pol}}(\Phi) - W_{\text{arith}}(\Phi) - W_{\text{arch}}(\Phi)
    \end{equation*}
    où les différentes fonctionnelles agissant sur $\Phi$ sont explicitées analytiquement par :
    \begin{align*}
        W_{\text{pol}}(\Phi) \;&=\; \mathcal{M}[\Phi](1) + \mathcal{M}[\Phi](0) \\
        W_{\text{arith}}(\Phi) \;&=\; \sum_{p \in \mathbb{P}} \sum_{n=1}^{\infty} \frac{\ln(p)}{p^{n/2}} \left( \Phi(p^n) + \Phi\left(\frac{1}{p^n}\right) \right) \\
        W_{\text{arch}}(\Phi) \;&=\; \left( \ln(\pi) + \gamma \right) \Phi(1) + \int_{1}^{\infty} \frac{\Phi(x) + \Phi\left(\frac{1}{x}\right) - 2\Phi(1)}{x - \frac{1}{x}} \, \frac{\mathrm{d}x}{x}
    \end{align*}
    avec $\mathcal{M}$ désignant la transformée de Mellin, $\mathbb{P}$ l'ensemble des nombres premiers et $\gamma$ la constante d'Euler-Mascheroni.
\end{DefNef}

À partir de cette définition, Weil a démontré que la validité de la conjecture de Riemann 
équivaut à la positivité d'une certaine forme fonctionnelle globale.

\begin{proposition}[label=pr:CriterePositiviteWeil_fr_intro]{Critère de positivité de Weil}
    Soit $\mathcal{Z}$ l'ensemble des zéros non triviaux de
    la fonction Zeta de Riemann.\par
    D'après le théorème d'André Weil~\cite{weil1952_formules}, 
    la forme hermitienne $\langle \cdot, \cdot \rangle_{\mathcal{HP}}$ 
    possède les propriétés suivantes : \par
    Pour toute fonction $f \in \mathcal{W}$, l'évaluation de cette forme sur la
    diagonale s'exprime analytiquement comme une somme sur l'ensemble
    $\mathcal{Z}$:
    \begin{equation*}
        \langle f, f \rangle_{\mathcal{HP}} \;=\; \sum_{\rho \in \mathcal{Z}} \mathcal{M}[f](\rho) \, \overline{\mathcal{M}[f](1 - \overline{\rho})}
    \end{equation*}
    où $\mathcal{M}[f]$ désigne la transformée de Mellin de la fonction test $f$. \par
    De plus, l'Hypothèse de Riemann est vérifiée si et seulement si, pour toute fonction 
    non nulle $f \in \mathcal{W}$, on a :
    \begin{equation*}
        \langle f, f \rangle_{\mathcal{HP}} \;\geq\; 0
    \end{equation*}
\end{proposition}

Ce résultat a fondamentalement transformé la nature du problème. 
Il a permis de traduire l'hypothèse de Riemann en une question d'analyse 
fonctionnelle et de théorie spectrale. Cette reformulation a inspiré de 
nombreuses tentatives de résolution basées sur la construction d'opérateurs 
dont le spectre correspondrait aux zéros de la fonction zêta, rejoignant 
ainsi la conjecture spectrale de Hilbert-Pólya.\medskip\par

Celle-ci suggère que les parties imaginaires des zéros non triviaux de la fonction 
zêta de Riemann correspondent aux valeurs propres d'un opérateur auto-adjoint 
(un hamiltonien) agissant sur un espace de Hilbert. Si un tel opérateur existe, 
la nature auto-adjointe garantit que ses valeurs propres sont réelles, ce qui 
implique immédiatement que tous les zéros se trouvent sur la droite critique, 
démontrant ainsi l'hypothèse de Riemann.

Par la suite, le mathématicien français Alain Connes a apporté une résolution 
à ce programme. En s'appuyant sur la géométrie non commutative, 
il a construit un espace adéquat et établi une formule de trace spectrale 
fournissant une réalisation géométrique et spectrale des zéros de la fonction zêta. 
Ses travaux ont ainsi donné un cadre concret à l'approche de Hilbert-Pólya.

Dans le présent document, nous allons reprendre en détail le développement des 
formules explicites et de la forme hermitienne d'André Weil. L'objectif principal 
est de dérouler l'intégralité du raisonnement menant à la démonstration de 
l'hypothèse de Riemann dans le cadre recherché. Pour ce faire, 
nous montrerons que le critère de positivité 
de Weil est effectivement vérifié pour un produit scalaire spécifique, dont nous 
donnerons la forme explicite.

Cette démarche s'articule notamment autour des deux résultats donnés 
par~\ref{lm:EquivalenceFormeWeil_fr} et~\ref{lm:ConstructionMetriqueHP_fr}, 
qui lient la distribution de Weil à la géométrie spectrale:

\begin{lemma}[label=lm:EquivalenceFormeWeil_fr_intro]{Action de la distribution de Weil et équivalence spectrale}
    Soient $f$ et $g$ deux fonctions de l'espace de Weil $\mathcal{W}$. \par
    Soit la distribution de Weil $\mathcal{K}_{\Gamma}$ définie sur le produit $\mathbb{R}_+^* \times \mathbb{R}_+^*$ par la somme de ses composantes (identité, arithmétique et archimédienne) :
    \begin{equation*}
        \mathcal{K}_{\Gamma}(x,y) \;=\; \delta(x-y) \;-\; \sum_{n=1}^{\infty} \frac{\Lambda(n)}{\sqrt{n}} \left[ \delta\left(x - ny\right) + \delta\left(y - nx\right) \right] \;+\; \mathcal{K}_{\infty}(x,y)
    \end{equation*}
    où $\Lambda$ est la fonction de von Mangoldt (Définition~\ref{df:VonMangoldt_fr}), 
    $\mathcal{K}_{\infty}$ le noyau archimédien (Définition~\ref{df:NoyauArchimedien_fr}), 
    et $\delta$ la distribution de Dirac 
    (Définition~\ref{df:DistributionDiracMultiplicative_fr}). \par
    On définit l'action de cette distribution sur le produit tensoriel $f \otimes \overline{g}$ par rapport à la mesure de Haar multiplicative invariante $\frac{\mathrm{d}x}{x}$ :
    \begin{equation*}
        \langle \mathcal{K}_{\Gamma}, f \otimes \overline{g} \rangle \;:=\; \int_{0}^{\infty} \int_{0}^{\infty} f(x) \overline{g(y)} \, \mathcal{K}_{\Gamma}(x,y) \, \frac{\mathrm{d}x}{x} \frac{\mathrm{d}y}{y}
    \end{equation*}
    D'après les travaux d'André Weil~\cite{weil1952_formules}, cette intégrale double converge absolument et vérifie l'égalité spectrale suivante sur l'ensemble $\mathcal{Z}$ des zéros non triviaux de la fonction zêta de Riemann :
    \begin{equation*}
        \langle \mathcal{K}_{\Gamma}, f \otimes \overline{g} \rangle \;=\; \sum_{\rho \in \mathcal{Z}} \mathcal{M}[f](\rho) \, \overline{\mathcal{M}[g](1 - \overline{\rho})}
    \end{equation*}
    Par conséquent, l'action de la distribution coïncide exactement avec la forme
    sesquilinéaire de Weil associée $\langle f, g \rangle_{\mathcal{HP}}$,
    décrite dans la
    proposition~\ref{pr:CriterePositiviteWeil_fr},
    et qui a été définie dans la définition~\ref{df:FormeHermitienneWeil_fr}.
\end{lemma}

\begin{lemma}[label=lm:ConstructionMetriqueHP_fr_intro]{Structure de Hilbert-Pólya et opérateur hamiltonien}
    Soit $\mathcal{Z}$ l'ensemble des zéros non triviaux de la fonction zêta de Riemann.\par
    D'après le lemme~\ref{lm:EquivalenceFormeWeil_fr},
    l'action de la distribution bilinéaire de Weil sur l'espace $\mathcal{W}$
    coïncide avec la forme sesquilinéaire spectrale, que l'on
    notera $\langle \cdot, \cdot \rangle_{\mathcal{HP}}$. \par
    D'après la proposition~\ref{pr:CriterePositiviteWeil_fr} (Critère de positivité
    de Weil), cette forme hermitienne vérifie $\langle f, f \rangle_{\mathcal{HP}}
    \geq 0$ pour toute fonction $f \in \mathcal{W}$ si et seulement si l'hypothèse
    de Riemann est vérifiée. \par
    Sous cette condition, soit $\mathcal{N}$ le noyau associé à cette forme
    (sous-espace isotrope). \par
    Comme la forme $\langle \cdot, \cdot \rangle_{\mathcal{HP}}$ induit un produit
    scalaire strictement défini positif sur l'espace quotient
    $\mathcal{W} / \mathcal{N}$, on en déduit que la complétion topologique de cet
    espace quotient engendre un espace de Hilbert séparable, que l'on notera
    $\mathcal{H}_{HP}$. \par
    Sur cet espace hilbertien $\mathcal{H}_{HP}$, il existe un opérateur
    hamiltonien spectral non borné, que l'on notera $\mathbf{H}$, dont les
    valeurs propres correspondent exactement aux parties imaginaires des zéros
    non triviaux de la fonction zêta. \par
    Cet opérateur $\mathbf{H}$ est essentiellement auto-adjoint.
\end{lemma}

L'étude de ce document s'achèvera par l'énoncé et la démonstration du résultat suivant, 
qui permettra d'établir l'hypothèse de Riemann dans le théorème~\ref{th:ResolutionConjectureHP_fr}:

\begin{theorem}[label=th:ResolutionConjectureHP_fr_intro]{Réalisation spectrale géométrique et démonstration de l'Hypothèse de Riemann}
    Soient $\mathcal{H}$ un espace de Hilbert séparable et $\mathbf{D}$ un opérateur non borné, auto-adjoint et à résolvante compacte sur $\mathcal{H}$. \par
    On suppose que l'opérateur $\mathbf{D}$ engendre une formule de trace vérifiant l'identité d'absorption spectrale suivante, pour toute fonction test $f \in \mathcal{W}$ :
    \begin{equation*}
        \mathrm{Tr}_{\mathrm{reg}}\left( f(\mathbf{D}) f^*(\mathbf{D}) \right) \;=\; \mathcal{K}_\Gamma(f * f^*) \;=\; \sum_{\rho \in \mathcal{Z}} \mathcal{M}[f](\rho) \overline{\mathcal{M}[f](1-\bar{\rho})}
    \end{equation*}
    où $f^*(x) = \overline{f(x^{-1})} x^{-1}$ est l'involution sur l'algèbre de Schwartz multiplicative, et $\mathrm{Tr}_{\mathrm{reg}}$ dénote la trace régularisée sur l'espace de Hilbert $\mathcal{H}$. \par
    Alors, la forme hermitienne $\langle \cdot, \cdot \rangle_{\mathcal{HP}}$ est positive ou nulle sur $\mathcal{W}$, et par conséquent, l'Hypothèse de Riemann est vraie.
\end{theorem}

Pour les hypothèses du théorème, on se placera dans un 
idéal de Dixmier de sorte que la trace régularisée ait les 
propriétés nécessaires à la réalisation de celle-ci.


\end{document}